\section{Harmonic helpers and character-subgroup units}
\label{sec:harmonic-helpers}

\Cref{sec:findings} and \cref{sec:character-clusters} established that the
trained model implements modular multiplication via the multiplicative
Fourier algorithm on $\Zpstar{p}$, with a sparse set $\mathcal{K}$ of key
characters. At wide capacity ($d_{\text{mlp}}$ large) each $k \in \mathcal{K}$
has a clean dedicated MLP cluster computing $\chi_k(ab)$, and these clusters
correspond to mutually orthogonal write directions in the residual stream.

This note documents an effect that appears under tight capacity. At
$d_{\text{mlp}} = 32$, the embedding spectrum exposes a set of five
essential characters, but only three of them function as standalone Fourier
products. The other two are \emph{harmonic helpers}: characters that, when
removed from $\WE$, degrade the model's output at a \emph{different}
character's frequency.  The relation is structural --- the helper is
algebraically the square of the primary it boosts --- and admits a clean
group-theoretic reading.  The model under capacity pressure does not use a
sparse set of individual characters; it uses a small number of
\emph{cyclic-character subgroups}, and the helper characters are subgroup
generators that appear alongside their primaries.

\subsection{Empirical setup}
\label{subsec:harmonic-setup}

We use the $d_{\text{model}} = 24$, $d_{\text{mlp}} = 32$, seed 1 model
trained for $10{,}000$ epochs to test loss $\approx 1.8 \times 10^{-4}$.
The W$_E$ spectrum gives $\mathcal{K} = \{3, 10, 51\}$ at a $5\%$-of-max
threshold, but a full per-character W$_E$-ablation sweep
(\cref{subsec:harmonic-ablation}) shows two more characters --- $k = 20$
and $k = 6$ --- are also essential, with comparable test-loss impact to
$k = 51$. We refer to the augmented set
$\widetilde{\mathcal{K}} = \{3, 10, 20, 51, 6\}$ throughout.

\subsection{Direct ablation of input characters}
\label{subsec:harmonic-ablation}

For each character $k$ we form an ablated embedding $\WE^{\setminus k}$ by
projecting $\WE$ onto the orthogonal complement of the $\cos_k$ and $\sin_k$
basis directions (i.e., zeroing out W$_E$'s projection onto character $k$).
Define
\begin{equation}
  \Delta_k(a, b, c)
  \;:=\; \mathrm{logit}\bigl(c\,\big|\,a, b\,;\, \WE\bigr)
  \;-\; \mathrm{logit}\bigl(c\,\big|\,a, b\,;\, \WE^{\setminus k}\bigr).
  \label{eq:delta-k-def}
\end{equation}
By construction $\Delta_k$ is the exact contribution of character $k$ in
$\WE$ to each logit, with no projection-direction assumptions.

Reducing $\Delta_k(a, b, c)$ over the input grid: define
$j_x := \dlog_g(x)$ and average $\Delta_k$ over all $(a, b, c)$ with the
same $\Delta j := (j_c - j_a - j_b) \bmod (p-1)$:
\begin{equation}
  \widehat{\Delta}_k(\Delta j)
  \;:=\; \mathop{\mathrm{mean}}_{j_a + j_b - j_c \equiv -\Delta j}
            \Delta_k(a, b, c).
  \label{eq:delta-k-reduction}
\end{equation}
If character $k$ contributes a clean Fourier product to the logits, then
$\widehat{\Delta}_k(\Delta j) \propto \cos(2\pi k \,\Delta j / (p - 1))$
--- a pure cosine at frequency $k$.

\Cref{tab:harmonic-ablation-spectrum} reports the dominant frequency of
$\widehat{\Delta}_k$ for each essential $k$.

\begin{center}
\begin{tabular}{r|r|r|l}
  ablated $k$ & order & dominant freq in $\widehat{\Delta}_k$ & interpretation \\
  \hline
  3   & 112 & 3   & primary; clean Fourier product \\
  10  & 56  & 10  & primary; clean Fourier product \\
  51  & 112 & 51  & primary; clean Fourier product \\
  20  & 28  & \textbf{10} & helper; destroys $k = 10$ at output \\
  6   & 56  & \textbf{3}  & helper; destroys $k = 3$ at output \\
\end{tabular}
\captionof{table}{Top Fourier coefficient of the 1-D reduction
$\widehat{\Delta}_k$ for each essential character. The piggyback
characters $k = 20$ and $k = 6$ destroy output content at characters
$k = 10$ and $k = 3$ respectively, not at their own frequencies.}
\label{tab:harmonic-ablation-spectrum}
\end{center}

The asymmetry is striking. Ablating each of the three dominant primaries
($3, 10, 51$) destroys output content at the corresponding character's
frequency. Ablating either helper destroys output content at the
\emph{primary's} frequency. In both cases the helper has exactly twice the
frequency of the primary it boosts: $20 = 2 \cdot 10$ and $6 = 2 \cdot 3$
(modulo $p - 1 = 112$).

\subsection{The character-group reading}
\label{subsec:harmonic-group}

Let $G = \Zpstar{p}$ and write $n := p - 1$. Via the discrete logarithm
$G \cong \Zp{n}$, so the character group $\widehat{G}$ also satisfies
$\widehat{G} \cong \Zp{n}$, with the bijection
$k \longleftrightarrow \chi_k$.  Characters compose multiplicatively as
functions, and on indices this is addition in $\Zp{n}$:
\begin{equation}
  \chi_{k_1} \cdot \chi_{k_2} = \chi_{k_1 + k_2}, \qquad
  \chi_k^{\,m} = \chi_{m k}.
  \label{eq:char-product}
\end{equation}
In particular $\chi_{2k} = \chi_k^{\,2}$. The two helper-primary pairs
observed at $d_{\text{mlp}} = 32$ each consist of a character and its
square in $\widehat{G}$:
\begin{equation}
  \chi_6 \;=\; \chi_3^{\,2}, \qquad \chi_{20} \;=\; \chi_{10}^{\,2}.
\end{equation}

The cyclic subgroup $\langle \chi_k \rangle \le \widehat{G}$ generated by
the primary character has order $\ord(k)$ and contains all powers
$\{\chi_0, \chi_k, \chi_{2k}, \chi_{3k}, \ldots, \chi_{(\ord(k)-1) k}\}$.
The helper $\chi_{2k} = \chi_k^{\,2}$ is the \emph{first non-identity
square} inside this subgroup.  It generates the unique index-2 subgroup of
$\langle \chi_k \rangle$ (when $\ord(k)$ is even):
\begin{equation}
  \langle \chi_k^{\,2} \rangle
  \;=\; \{\chi_0, \chi_{2k}, \chi_{4k}, \ldots\}
  \;\le\; \langle \chi_k \rangle,
  \qquad [\langle \chi_k \rangle : \langle \chi_k^{\,2} \rangle] = 2.
\end{equation}

So the observed harmonic-helper pairs are exactly the \emph{first step of
the Sylow-2 filtration} of the cyclic subgroup generated by the primary.
At $p = 113$, where $n = 112 = 2^4 \cdot 7$, the Sylow-2 part of
$\widehat{G}$ has structure $\Zp{16}$, and the powers
$\chi_k \to \chi_{2k} \to \chi_{4k} \to \chi_{8k} \to \chi_{16k}$ walk
through this 2-part, halving the order at each step.

\paragraph{Dual interpretation.}
Pontryagin duality identifies $\widehat{G}$ with $G$ for finite abelian
groups, with cyclic subgroups of $\widehat{G}$ dual to quotients of $G$
(equivalently, to subgroups of $G$ under self-duality for cyclic groups).
The index-2 subgroup $\langle \chi_k^{\,2} \rangle$ of $\langle \chi_k
\rangle$ is dual to the subgroup of squares $G^2 \le G$, which for
$G = \Zpstar{p}$ is exactly the group of \emph{quadratic residues}. So
the helper character $\chi_{2k}$ is the restriction of $\chi_k$ to the
quadratic-residue subgroup --- the part of $\chi_k$ that is invariant
under multiplication by quadratic non-residues. Adding the helper to
$\WE$ provides the model with a direct algebraic handle on the
QR-restricted component of the primary character.

\subsection{Why the helper helps: squared features and polarization}
\label{subsec:harmonic-mechanism}

Write $\varphi_x := 2\pi\, j_x / (p-1)$ for the discrete-log phase of
$x \in G$, so that $\chi_k(x) = e^{i k \varphi_x}$ and the real basis
functions are $\cos(k\varphi_x), \sin(k\varphi_x)$. The frequency-$k$
contribution the model must place in the logits is the character product
\begin{equation}
  M_k(a,b) \;=\; \cos\bigl(k(\varphi_a + \varphi_b)\bigr),
  \label{eq:char-product-target}
\end{equation}
which the unembedding reads against $\cos(k\varphi_c), \sin(k\varphi_c)$ to
peak at $\varphi_c = \varphi_a + \varphi_b$, i.e.\ $c = ab$. Because
\eqref{eq:char-product-target} is a \emph{product} of an $a$-feature and a
$b$-feature, no linear map can form it; it must come from the MLP
nonlinearity.

\paragraph{The doubled character is the squared feature.}
The leading nonlinear action of a ReLU is quadratic. Expanding a single
rectified cosine,
\begin{equation}
  \mathrm{ReLU}(\cos\theta)
  \;=\; \frac{1}{\pi} + \frac{1}{2}\cos\theta
       + \frac{2}{3\pi}\cos 2\theta
       - \frac{2}{15\pi}\cos 4\theta + \cdots,
  \label{eq:relu-fourier}
\end{equation}
the first generated harmonic sits at $2\theta$ with coefficient
$\tfrac{2}{3\pi} \approx 0.212$ for a single unit (an antisymmetric pair
$\mathrm{ReLU}(z) + \mathrm{ReLU}(-z) = |z|$ doubles this to
$\tfrac{4}{3\pi} \approx 0.424$ and removes the fundamental).%
\footnote{An earlier draft quoted $\tfrac{4}{3\pi}$ for a single ReLU; the
correct single-unit value is $\tfrac{2}{3\pi}$, verified by dense FFT of
$\mathrm{ReLU}(\cos\theta)$.}
Squaring a character doubles its index, $\chi_k^{\,2} = \chi_{2k}$, and the
two statements are the same identity:
\begin{equation}
  \cos 2k\varphi \;=\; 2\cos^2 k\varphi - 1
  \quad\Longleftrightarrow\quad
  \cos^2 k\varphi \;=\; \tfrac{1}{2}\bigl(1 + \cos 2k\varphi\bigr).
  \label{eq:double-angle}
\end{equation}
So the helper $\chi_{2k}$ is precisely the \emph{squared} primary feature:
having $\chi_{2k}$ in $\WE$ makes $\cos^2(k\varphi)$ available as a
\emph{linear} readout.

\paragraph{Products are built from squares.}
With $x = \cos k\varphi_a$, $y = \cos k\varphi_b$, $u = \sin k\varphi_a$,
$v = \sin k\varphi_b$, the target \eqref{eq:char-product-target} is
\begin{equation}
  M_k(a,b)
  = \cos k\varphi_a \cos k\varphi_b - \sin k\varphi_a \sin k\varphi_b
  = xy - uv.
  \label{eq:product-xy}
\end{equation}
Products decompose into squares by polarization,
\begin{equation}
  xy = \tfrac{1}{2}\bigl[(x+y)^2 - x^2 - y^2\bigr],
  \label{eq:polarization}
\end{equation}
and the elementary nonlinear operation a one-hidden-layer ReLU MLP performs
is approximating a square: a width-$W$ ReLU layer fits a parabola on a
bounded interval to error $\sim W^{-2}$, i.e.\ needs $\sim
\varepsilon^{-1/2}$ units for error $\varepsilon$.

\paragraph{What the helper buys.}
By \eqref{eq:double-angle} the self-squares $x^2, y^2, u^2, v^2$ are all
\emph{affine} in the helper features $\cos 2k\varphi_a, \cos 2k\varphi_b$.
With $\chi_{2k}$ present in $\WE$ the MLP therefore obtains four of the six
squares needed by \eqref{eq:polarization} (counting the $\sin$ partner
$uv$) for free, and must approximate only the two cross-squares $(x+y)^2$
and $(u+v)^2$ with ReLU units. Without $\chi_{2k}$ it must synthesise all
six. The helper thus removes a constant factor (here $6 \to 2$) of the
neurons needed to build a clean $\chi_k$ product --- a saving that bites
hardest exactly when neurons are scarce.

\paragraph{Numerical confirmation.}
\Cref{fig:helper-fvu} fits a width-$W$ ReLU MLP to the bare product
\eqref{eq:char-product-target} for $k = 10$ on the full group grid, with
input features $\{\chi_{10}\}$ versus $\{\chi_{10}, \chi_{20}\}$. The helper
cuts the fraction of variance unexplained by roughly $10\times$ at width $3$
($0.39 \to 0.038$); the gap is largest at small width and closes as $W$
grows (both below $1\%$ by $W = 8$). This is the capacity-pressure signature
in a setting with no trained model: when neurons are scarce the precomputed
square is decisive, and when they are plentiful the MLP builds the square
itself and the helper becomes redundant.

\begin{center}
\includegraphics[width=0.78\linewidth]{helper_mechanism_fvu}
\captionof{figure}{Approximating the character product
$\cos(10(\varphi_a+\varphi_b))$ with a width-$W$ ReLU MLP, given the primary
character alone (red) versus the primary plus its helper $\chi_{20}$ (blue).
FVU is mean-squared error normalised by target variance ($= 1 - R^2$). The
helper's advantage is concentrated at small width.}
\label{fig:helper-fvu}
\end{center}

\paragraph{Consistency with the ablation.}
This is why ablating $\chi_{2k}$ from $\WE$ (\cref{eq:delta-k-def}) destroys
output content at frequency $k$ rather than $2k$: removing the helper
removes the precomputed squared features, and a finite-width MLP can no
longer assemble a clean $\chi_k$ product. The dependence is genuinely
cross-frequency and nonlinear --- an additive, single-frequency model in
which each neuron contributes independently at one frequency \emph{cannot}
express it. The harmonic helper is therefore a signature of the ReLU
nonlinearity coupling a character to its square, not an independent Fourier
mode.

\paragraph{The squaring ladder.}
Iterating, $\chi_k \to \chi_{2k} \to \chi_{4k} \to \chi_{8k} \to \cdots$ is
repeated squaring: it supplies $\cos^2, \cos^4, \ldots$ of the primary
feature, the higher-degree building blocks a narrower MLP needs to reach a
given product accuracy. Two consequences follow. First, the chain the
nonlinearity \emph{naturally} generates is the doubling (Sylow-2) ladder,
because squaring multiplies the character index by $2$; a Sylow-$7$ step
($\times 7$) would require a degree-$7$ term, which a quadratic
nonlinearity does not produce at meaningful amplitude. Second, tighter
capacity recruits \emph{deeper} into the ladder --- which is exactly the
subgroup-unit picture of \cref{subsec:harmonic-subgroup-unit}.

\subsection{The subgroup-unit picture}
\label{subsec:harmonic-subgroup-unit}

Putting the empirical and group-theoretic pieces together: at
$d_{\text{mlp}} = 32$ the model does not use a sparse set of individual
Fourier characters. It uses a small collection of \emph{cyclic-character
subgroups} as algorithmic units, each generated by a primary character
and augmented by helpers that walk down the Sylow-2 filtration:
\begin{center}
\begin{tabular}{c|c|c}
  primary $\chi_k$ & helpers (subgroup steps) & terminates at \\
  \hline
  $\chi_3$  (order 112) & $\chi_6$ (order 56)  & {\small possibly $\chi_{12}, \chi_{24}, \ldots$ at smaller $d_{\text{mlp}}$} \\
  $\chi_{10}$ (order 56) & $\chi_{20}$ (order 28) & {\small possibly $\chi_{40}, \ldots$} \\
  $\chi_{51}$ (order 112) & --- & {\small (none observed; capacity exhausted)} \\
\end{tabular}
\end{center}

The number of subgroup generators stored in $\WE$ for each primary is a
function of the d\_mlp budget. Wide capacity stores only the primary;
tight capacity stores deeper into the Sylow-2 chain.

This reframes \cref{sec:findings} item 3 (`$\mathcal{K}$ is sparse'):
$\mathcal{K}$ is sparse \emph{in characters}, but its structure under
capacity pressure is better described as sparse \emph{in cyclic
subgroups}, with characters drawn as needed from each chosen subgroup's
Sylow-2 chain.

\subsection{Cross-budget verification and the Sylow-2 specificity}
\label{subsec:harmonic-verification}

We tested the same $\Delta_k$ ablation at two other architecture settings.

At $d_{\text{mlp}} = 128$, seed 1, with $\mathcal{K} = \{5, 10, 12, 49,
52\}$: the predicted pair $(5, 10)$ (since $10 = 2 \cdot 5$ and both
have non-trivial Sylow-2 component) appears. Ablating $k = 10$ destroys
output content at frequency $5$, exactly as the framework predicts. The
magnitude of the helper effect is small ($E \approx 2 \times 10^2$,
compared to primary effects of $10^5$--$10^6$), consistent with the
hypothesis that helpers diminish under abundant capacity.

At $d_{\text{mlp}} = 64$, seed 1, with $\mathcal{K} = \{5, 12, 32, 56\}$:
each primary destroys its own frequency cleanly at magnitude $10^5$--$10^6$,
and there are no detectable helper effects. The model uses four
independent primaries with no shared subgroup chains.

These results refine the claim of \cref{subsec:harmonic-subgroup-unit}:
helpers appear when capacity is tight relative to the number of primaries,
and only along the Sylow-2 chains. The character $k = 32$ at
$d_{\text{mlp}} = 64$ is purely in the Sylow-7 part of $\widehat{G}$
(order 7); doubling in $\widehat{G}$ along the Sylow-7 direction moves
within the same order-7 subgroup, so no proper subchain exists to walk
down. Purely Sylow-7 primaries do not admit Sylow-2-style helpers, which
is consistent with the absence of any helper for $k = 32$. Mechanistically
(\cref{subsec:harmonic-mechanism}) this is the same statement as the
squaring ladder: ReLU's quadratic action generates index-doublings, not
multiplications by $7$, so the nonlinearity can build a Sylow-2 chain but
not a Sylow-7 one.

\subsection{Open questions}
\label{subsec:harmonic-open}

\begin{itemize}
\item \textbf{The capacity threshold.} The empirical pattern across
  $d_{\text{mlp}} \in \{32, 64, 128\}$ shows non-monotone helper usage:
  helpers present and substantial at $32$, absent at $64$, present and
  weak at $128$. The likely explanation is the interaction between
  $|\mathcal{K}|$ and per-character neuron budget: when $|\mathcal{K}|$
  is large enough that per-primary neuron count falls below the
  approximation floor, helpers are recruited. Quantifying this
  threshold requires sweeping $d_{\text{mlp}}$ at fixed $|\mathcal{K}|$
  (or vice versa) and measuring helper amplitude as a function of
  $d_{\text{mlp}} / |\mathcal{K}|$.

\item \textbf{Sylow-7 helpers.} The Sylow decomposition of
  $\widehat{G} = \Zp{112}$ is $\Zp{16} \times \Zp{7}$. The Sylow-7 part
  supports the chain $\chi_k \to \chi_{7k}$ (multiplying index by $7$
  mod $112$). We have only observed the Sylow-2 step in our experiments
  to date; whether the Sylow-7 chain also produces helpers is open.

\item \textbf{Connection to capacity floor.} If the model uses
  subgroup-chains as the algorithmic unit, the d\_mlp lower bound for
  grokking depends on how many characters per chain the architecture
  must store, not on the number of distinct primaries. A primary
  with deep Sylow-2 chain might require more $\WE$ slots but fewer MLP
  neurons per character than a primary with no helpers. The interaction
  between primary diversity and chain depth is the right frame for the
  empirical d\_mlp floor.

\item \textbf{Non-abelian generalization.} For non-abelian groups, the
  analog of cyclic subgroups in $\widehat{G}$ is replaced by the lattice
  of subgroups in the set of irreducible representations. Whether SGD
  finds chains of irreducibles closed under tensor square --- the
  natural generalization of $\chi_k \mapsto \chi_k^{\,2}$ ---
  remains untested.
\end{itemize}
